\documentclass[10pt]{article}
\usepackage[margin=20mm]{geometry}
\usepackage{amsmath,amssymb,mathtools,fontspec,unicode-math,xurl,hyperref}
\IfFontExistsTF{DejaVu Serif}{\setmainfont{DejaVu Serif}}{\setmainfont{Latin Modern Roman}}
\setmathfont{latinmodern-math.otf}
\setlength{\parindent}{0pt}\setlength{\parskip}{6pt}\sloppy
\title{Technical Verification of
Mapping Class Group of Manifolds Which Look Like
3-Dimensional Complete Intersections}
\author{}
\date{}
\begin{document}
\maketitle
Abstract: The acquired author manuscript contains a consequential arithmetic error in its nonspin kernel table and incompatible generator relations in its presentation theorem. An admissible complete intersection of degrees \((2,3)\) gives an exact order-\(21\) characteristic-number class where the printed table gives order \(3\). The quintic likewise exposes the incorrect presentation orders and missing relations. Two impossible relative dual-class identifications occur in the geometric calculation underlying the kernel description. Local corrections to a characteristic number, a genus-zero splitting clause, and an integral symplectic homology table are also established. Counterevidence and the limits of proposed repairs are retained. Scientific adjudication remains pending.\par
\section*{Paper identity and inspected scope}
The paper is by Matthias Kreck and Yang Su, with journal identifier 10.1215/00127094-2024-0036 (https://doi.org/10.1215/00127094-2024-0036). The scientific object inspected is [KS], arXiv:2009.08054v3, 11 June 2024, 62 pages. All page references below are its one-based PDF pages, which agree with its printed numbering. The primary PDF has 576,226 bytes. Its matching source archive and extracted `source.tex` were available and verified. The manuscript’s description as a final version does not establish byte-level equivalence with the publisher’s article. The publisher endpoint did not yield an independently inspected article, so the findings are confined to this author version.\par
The complete manuscript text was read. Detailed checks addressed the variation lattice and extension, characteristic-number conventions, the spin and nonspin relation lattices, explicit generators, the relative-to-absolute kernel argument, low-genus boundaries, lower homology, and the bordism appendix. Exact integer calculations supplement the arguments below. Selected primary dependencies were acquired and inspected at the stated locations; a complete reconstruction of all cited surgery, embedding, and arithmetic-group results is not claimed. Source inventory parsing and hashing are separate from scientific checking.\par
\section*{Definitions and operative assumptions}
The standing class consists of smooth, oriented, closed, simply connected six-manifolds \(M\) with \(H_2(M;\mathbb Z)\cong\mathbb Z\), a chosen generator \(x\in H^2(M;\mathbb Z)\), and \(d=\langle x^3,[M]\rangle>0\). Write \(x^2=dy\), with \(y\) the corresponding generator of \(H^4(M;\mathbb Z)\), and \(H=H^3(M;\mathbb Z)\), a free symplectic lattice of rank \(2g=b_3(M)\). In the spin case, \[p_1(TM)=4ky,\qquad k=d+6l.\] In the nonspin case, \[d=2\delta,\qquad p_1(TM)=(4k-d)y=(d+48l)y,\qquad
 k=\delta+12l.\] Here \(l\in\mathbb Z\). The Torelli group \(\mathcal I(M)\) is the kernel of the action on \(H\). The variation takes values in \[L_{d,k}=(H\oplus H)/\Delta_{d,k},\qquad
 \Delta_{d,k}=\{(du,ku):u\in H\},\] and \(\mathcal K(M)\) is its kernel. The special Torelli group \(\mathcal J(M)\) acts trivially on \(\pi_3(M)\) and fits into \(1\to\mathcal K(M)\to\mathcal J(M)\to H\to0\). A fixed disk gives the relative group \(\mathcal K(M,D)\).\par
Normal and tangent Pontryagin classes are distinguished throughout. In particular, for a spin bordism, \(p_1(\nu W)=-p_1(TW)\). The relative characteristic numbers require actual preimages of the specified absolute classes. Replacing such a preimage by a class with different absolute evaluations is not an allowable choice ambiguity.\par
\section*{Substantive findings}
Problem 1 (Major: the nonspin table loses a factor of seven).  Theorem 6.1(2), pp. 26–27, and its reduction in Section 10.4, pp. 51–52, disagree with the original geometric relation lattice (10.5), p. 49, on a smooth complete intersection satisfying every standing hypothesis.\par
\subsection*{Decisive calculation}
Let \(M=X_{2,3}\subset\mathbb {CP}^5\) be a smooth complete intersection of degrees \(2\) and \(3\). It is a simply connected complex threefold with \(H_2(M;\mathbb Z)=\mathbb Z\). The hyperplane class is primitive, and the tangent/normal sequence gives \[\begin{gathered}
 d=6,\qquad c_1(TM)=(6-2-3)x=x,\\
 p_1(TM)=(6-2^2-3^2)x^2=-7x^2=-42y .
 \end{gathered}\] Thus \(w_2(M)\ne0\), \(\delta=3\), \(k=-9\), and \(l=-1\). This is an actual complete intersection, not a formal parameter pair.\par
Equation reference: eq:21. Put \(h=d+24l\). The columns of (10.5) are \[\begin{aligned}
 c_1&=(0,0,-224),&c_2&=(6,-6,8),&c_3&=(1,-3,8),\\
 c_4&=(d^2,-d(h-1),(h-1)^2-1),\\
 c_5&=(2d^2,-2dh,2h^2),&
 c_6&=(-2d,3d+24l,-4h).
\end{aligned}\] For this \(M\), the last three columns are \[(36,114,360),\quad (72,216,648),\quad (-12,-6,72).\] The geometric generator \(N_{-1}\) in Section 10.4 has original characteristic vector \(v=(1,-1,0)\): its \((\mu/4,\lambda)\) coordinates are \((-1,-2)\), as on p. 51. Its order modulo the original column lattice is exactly \(21\). Indeed, \[
 21v=9c_3-c_6.\] Every column has second coordinate divisible by \(3\), whereas \(v_2=-1\). Also the functional \[\phi(\sigma_1,\sigma_2,\sigma_3)=2\sigma_1+\sigma_2+\sigma_3\pmod7\] annihilates every column and satisfies \(\phi(v)=1\). Hence the order is divisible by both \(3\) and \(7\), and [eq:21] (\#eq:21) gives the matching upper bound. Exact determinant computations independently give indices \(336\) and \(16\) before and after adjoining \(v\).\par
The source of the discrepancy is directly visible. Applying the printed definition \(\mu=6\sigma_1+10\sigma_2+3\sigma_3\) to \(c_5\) gives \[\frac{\mu(c_5)}4=-2(\delta^2+48\delta l-432l^2).\] The transformed table on p. 51 instead has \(-2(\delta^2+96\delta l-864l^2)\). At \((\delta,l)=(3,-1)\) these values are \(1134\) and \(2286\), respectively. The latter incorrectly kills the seven-primary class. Theorem 6.1 assigns \(a=0,b=1,c=0\), hence \(\mathcal K(M)\cong\mathbb Z/3\).\par
\subsection*{Geometric verification, impact, and repair}
The original columns were checked from \(\alpha=x^2-dy\) and \(\beta=(p_1(\nu W)-x^2)/2+hy\). For example, on \(\mathbb {CP}^2\times S^4\), \(\beta=-2x^2+hy\), giving \((-2d,3d+24l,-4h)\). On \(S^4\times S^4\), \(x=0\), \(y=a+b\), and \(\langle y^2,[W]\rangle=2\), giving \(c_5\). On \(\mathbb {HP}^2\) with its positive generator \(y\), \(\beta=(h-1)y\); subtracting the constant-map \(\mathbb {HP}^2\) gives \(c_4\). The \(\mathbb {CP}^2\times\mathbb {CP}^2\) and \(\mathbb {CP}^4\) calculations give \(c_2,c_3\). The nonspin sphere-bundle construction at \(m=-1\) gives \((1,-1,0)\) directly. Thus the discrepancy is not merely a comparison of two unmotivated tables.\par
The invariant homomorphism and the gluing comparison of Proposition 10.2 and Lemma 10.4 identify this class with the cyclic kernel generator. Within that framework, the answer is \(21\), not \(3\). At minimum the original lattice, its claimed simplification, and the theorem cannot all be correct; the explicit arithmetic and geometric column calculations locate the error in the simplification. The full external bordism-generation and modified-surgery theorems have not been reproved from first principles.\par
The corrected seven-primary condition from this lattice is \(\delta\equiv4l\pmod7\), equivalently \(d\equiv l\pmod7\), rather than \(d\equiv l\equiv0\pmod7\). The \(c_6\) coordinate is \(3(\delta-4l)\); when it vanishes modulo \(7\), all other relevant \(\mu/4\) coordinates vanish as well. A separate discrepancy in the \(\lambda(c_4)\) coefficient is a multiple of \(4\) and is harmless under the fixed \(\lambda\) relation. The quadric parameters \((d,l)=(2,0)\) give the same trivial kernel in both calculations. Repair requires correcting the transformation, the seven-primary statement, and its dependent applications; it does not require discarding every entry of the table.\par
Problem 2 (Major: Theorem 2.16 does not present the stated group).  The generator orders in Theorem 2.16, p. 10, repeated on p. 43, are incompatible with (8.1) and the paper’s quintic result. Changing only individual orders would still omit necessary joint relations.\par
For the quintic \(X_5\), the verified parameters are \((d,k,l)=(5,-25,-5)\), and (8.1) becomes the column lattice of \[C=\begin{pmatrix}
 -28&0&0&325&625&0\\
 0&12&0&140&275&-30\\
 0&0&2&-255&-500&50
 \end{pmatrix}.\] Its index is \(2\), with Hermite basis \(\left(\begin{smallmatrix}2&1&1\\0&1&0\\0&0&1\end{smallmatrix}\right)\). The character \((a,b,c)\mapsto a+b+c\pmod2\) kills every column and is nonzero on each coordinate vector. Each coordinate therefore represents the same nonzero element of \(\mathbb Z/2\).\par
The classes \(f_c,u,z\) are identified on p. 43 with these three coordinate vectors. Their orders are consequently \((2,2,2)\), whereas the printed formulas give \((7,6,1)\). This also contradicts Theorem 2.5, p. 4, which explicitly makes the disk-supported class nontrivial in the quintic’s \(\mathbb Z/2\) center. For an exact upper certificate, if \(e_i\) denotes the \(i\)th coordinate vector, then \[2e_1=C(91,-95,1035,4,2,-1)^t,\qquad
 2e_2=C(25,-24,265,6,-2,0)^t,\qquad 2e_3=c_3.\] These equalities and the parity character avoid reliance on numerical linear algebra.\par
The independent boundary example \(\mathbb {CP}^3\), with \((d,k,l)=(1,1,0)\), gives coordinate orders \((4,4,1)\) rather than \((7,12,1)\). The nonspin quadric gives a trivial kernel, whereas the nonspin order formula in Theorem 2.16 gives \(168\). For these genus-zero examples \(\mathcal J=\mathcal K\), so the contradiction cannot be attributed to an infinite quotient of \(\mathcal J\).\par
This is a defect in an explicit principal conclusion. Even substituting the correct quintic orders leaves a presentation with extra independent central generators: the actual relations include \(f_c=u=z\). A valid repair must recover the full relation lattice among the central generators and its interaction with the commutators, not just swap exponents or replace three orders. The abstract spin table agrees with the tested relation lattices; the incorrect presentation does not itself disprove that abstract table or the quintic description.\par
Problem 3 (Major: impossible relative dual classes in core calculations).  Section 7.3, pp. 36–37, and the proof of Proposition 8.1, pp. 39–40, use relative classes whose stated absolute images are incompatible with their intersection matrices. These steps feed the twist invariants and the surjectivity in Lemma 8.4 used to determine the kernel.\par
On p. 36 let \(e=[S_b^4]\) be the actual zero-section sphere and \(f=[Y]\). Their intersection form is \(Q=\left(\begin{smallmatrix}\chi&1\\1&0\end{smallmatrix}\right)\). Page 37 simultaneously requires \(j^*\bar\alpha=x^2\) and \(D\bar\alpha=e\). These assertions imply \[\langle x^2,e\rangle=e\cdot e=\chi .\] But \(H^2(S^4;\mathbb Z)=0\), so the restriction of \(x\) to the zero-section vanishes and \(\langle x^2,e\rangle=0\). The permitted twist parameter \(\chi=1\), used for \(\epsilon\), is decisive. This contradiction is independent of the tangent-versus-normal sign convention.\par
The same difficulty recurs on p. 40. The first two basis elements are the two zero-section spheres; the normal map has zero \(y\) evaluation on them. Thus \(\alpha=x^2-d\bar\nu^*y\) evaluates to zero on each. Nevertheless the proposed dual \(S_{\alpha_1}+S_{\alpha_2}\) has evaluations \[\begin{pmatrix}d&0&1&0\\0&-d&-1&0\\1&-1&0&0\\0&0&0&0\end{pmatrix}
 \begin{pmatrix}1\\1\\0\\0\end{pmatrix}
 =\begin{pmatrix}d\\-d\\0\\0\end{pmatrix}.\] Since \(d>0\), it cannot be a preimage of the specified \(\alpha\). The subsequent vanishing of characteristic numbers cannot be justified with that class.\par
There is useful repair evidence, but it is not a completed replacement of the geometry. In Section 7, if the Hopf evaluation is normalized by \(\langle x^2,Y\rangle=1\), the compatible dual is \(e-\chi f\). At the chosen framing \(A=0\), the stated normal class has dual \((2b-\chi)f\). Hence \[\bar\alpha^2=-\chi,\qquad
 \bar\alpha\bar\beta=2b-\chi,\qquad s_2/2=b-\chi .\] The last value is the negative of the printed one, but \(\rho\) still generates the cyclic order-\(12\) factor and \(\epsilon\) still has its order-\(2\) coordinate. The corresponding nonspin endpoint agrees with \(\mu/4=5b-4\chi\) on p. 52.\par
Two associated p. 37 inferences also require correction. Multiplication by the displayed matrix gives \(\bar\beta^2=A((2-A)\chi-4b)\), the opposite of the printed sign. Divisibility by \(8\) for all \(\chi,b\) forces \(A\) to be even, not \(A\equiv2\pmod4\): \(A=0\) is a direct counterexample to that inference. Evenness suffices for the stated reframing. For p. 40, the relative-image condition and the displayed radical force \(\alpha=0\) at the algebraic level. Moreover the spin characteristic property of \(p_1/2\) forces \(A\) even when \(Y^2=0\), so the proposed normalization \(A\in\{0,1\}\) would leave only \(A=0\). These observations suggest a shorter corrected calculation. They do not certify the required attaching maps, framing freedom, or normal extension. A new proof of these steps and propagation through Lemma 8.4 is needed. No additional counterexample to the final twist-generator theorem is asserted.\par
Problem 4 (Minor: the sphere-bundle invariant in (5.4)).  For the explicitly constructed \(N_m=S(\xi_m)\) on pp. 25–26, equation (5.4) omits \(8\) from \(s_1(N_m)\).\par
Here \(e(\xi_m)=x^2\), \(p_1(\xi_m)=(3+4m)x^2\), and \(p_1(T\mathbb {CP}^2)=3x^2\). Thus on the disk bundle \(p_1(\nu W)/2=-(2m+3)x^2\), and the Thom self-pairing is \(1\). The prescribed connected sum with \(\overline{\mathbb {HP}^2}\) makes the signature zero and subtracts \(1\) from the Pontryagin square. Therefore \[s_1(N_m)=(2m+3)^2-1=4m(m+3)+8,
 \quad s_2(N_m)=-2(m+1),\quad s_3(N_m)=-(2m+3).\] At \(m=0\), \(s_1=8\), rather than \(0\), in \(\mathbb Z/224\). The sign difference in \(s_3\) disappears modulo \(2\). The \(8\) difference is genuine for the unreduced invariant, but disappears modulo \(8\) after quotienting by homotopy spheres. Lemma 7.8 on p. 35 uses precisely that quotient. Thus this correction does not invalidate its reduced generation argument or, by itself, Theorem 7.9.\par
Problem 5 (Minor: the genus-zero nonsplitting boundary).  The residual nonsplitting clause of Theorem 4.6, p. 21, includes \(b_3(M)=0\), although the extension then splits.\par
The standing assumptions allow \(\mathbb {CP}^3\), for which \(H=0\). The Torelli group is the whole mapping class group, the variation target vanishes, and \(\mathcal K=\mathcal I=\operatorname{MCG}\). Both the kernel and quotient in (4.4) are trivial, so the extension splits. This needs no classification of \(\operatorname{MCG}(\mathbb {CP}^3)\). Include genus zero among the split cases and restrict the remaining nonsplitting cases to \(b_3\ge4\), as the proof’s actual case division does. Theorem 2.4 itself makes no contrary nonsplitting assertion at genus zero.\par
Problem 6 (Minor: missing genus-three integral torsion).  The table for \(E^2_{2,0}\) in Section 11.2, p. 55, incorrectly gives \(H_2(\operatorname{Sp}(2g,\mathbb Z);\mathbb Z)=\mathbb Z\) for every \(g\ge3\).\par
The directly cited primary [Retained latex: \textbackslash{}cite[Lemma 6.11 and the following table,
pp.\textasciitilde{}14--15]\{BCRR\}] gives \[H_2(\operatorname{Sp}(6,\mathbb Z);\mathbb Z)\cong\mathbb Z\oplus\mathbb Z/2,
 \qquad H_2(\operatorname{Sp}(2g,\mathbb Z);\mathbb Z)\cong\mathbb Z\quad(g\ge4).\] The genus-three summand must be retained in an integral spectral sequence. The later naturality argument factors the differential through zero coinvariants, so it kills this torsion as well. Rational rank is unchanged. This local correction does not itself refute the asserted integral surjectivity or rational dimension.\par
\section*{Unresolved scope and affirmative checks}
\subsection*{Corollary 2.10: Needs Evidence}
Corollary 2.10, p. 7, states \(\dim H_2(B\operatorname{Diff}_0(M);\mathbb Q)\ge2\) without an additional restriction. Its proof on p. 57 uses dimensions \(2\) for \(H_2(B\operatorname{MCG};\mathbb Q)\) and \(4\) for \(H_2(B\operatorname{Diff};\mathbb Q)\). The former is established here only for \(b_3\ge6\). The cited [Retained latex: \textbackslash{}cite[Section 5.3, pp.\textasciitilde{}19, 23--24]\{GRW\}] treats hypersurfaces \(V_d\subset\mathbb {CP}^4\) in degrees \[*\le (d^4-5d^3+10d^2-10d+4)/4=g(V_d)/2 .\] Its degree-two calculation supplies the needed input for these hypersurfaces with \(d\ge3\), not automatically for arbitrary standing \(M\) or low genus. Given the two input dimensions, the homology spectral sequence does yield the lower bound: the kernel in degree two is a quotient of fiber coinvariants. The issue is the scope of the inputs. An extension of the cited result or an additional argument is needed. No counterexample to the unrestricted corollary is asserted.\par
\subsection*{Checks that constrain the conclusions}
The exact sequence \(0\to H\to L_{d,k}\to H/d\to0\) and the injective doubling map \([a,b]\mapsto[a,2b]\) were checked directly, including possible torsion in \(L_{d,k}\). The quintic section \(v\mapsto[v,-5v]\) is well defined. The right-cocycle convention for the quadratic refinement was verified algebraically and by exact finite-field checks; this does not replace the external integral extension-class theorem.\par
The image-parity step has a direct supporting argument. For an oriented closed seven-manifold, instability gives \(v_4=0\), hence \(w_4=w_2^2\). For a spin mapping torus, \(q_1=p_1/2\) is therefore even; in the nonspin case the same holds for the spin bundle \(TM_f\oplus\eta\). The Wang sequence then makes the second variation even for every Torelli element. An alleged image gap is consequently not retained.\par
The intermediate center calculation on p. 57 omits \(d=1\), where \((H/d)\rtimes\operatorname{Sp}\) is \(\operatorname{Sp}\) and has central \(-I\). The paper’s next argument repairs this case: a lift of \(-I\) has commutator \(-2b\) with \(b\) in the free \(H\) kernel and cannot be central. The final center assertion is not disproved by this omission. The standard-module statements used here were checked in [Retained latex: \textbackslash{}cite[Appendix A, pp.\textasciitilde{}32--34]\{Krannich\}].\par
The total-degree-seven terms of the untwisted bordism spectral sequence give the stated upper bound \(24\). In the twisted sequence, the apparent objection from \(Sq^2b+xb=xb\) is invalid: \(xb=Sq^1(xa)\) is a Bockstein class and annihilates reduction of integral \(H_8\). The integral-row reduction in the differential matters; the relevant primary rule was checked in [Retained latex: \textbackslash{}cite[pp.\textasciitilde{}750--751]\{Teichner\}]. No counterexample to the appendix’s splitting was established.\par
\section*{Conclusion}
The nonspin kernel calculation and the presentation theorem require substantive correction. The explicit dual-class calculations also require replacement before their dependent geometric arguments can be accepted as written. Several broader constructions survive the particular tests reported here; this report does not infer wholesale failure from local errors.\par
Dimension | Recorded state\par
Initial assessment | Major issues established in the acquired author version\par
Repair verification | Exact arithmetic and local repairs checked; full revised geometric proofs and presentation not verified\par
Author adoption | Unknown; no contact or adoption claimed\par
Remaining evidence | Unrestricted Corollary 2.10 and unrebuilt external surgery, embedding, and arithmetic-group dependencies\par
Scientific adjudication | Pending; no independent acceptance\par
Human reviews | None recorded\par
\section*{Supporting checks}
The primary SHA256 is\par
> 00000TOKEN.\par
Failed attempts and corrections are retained. No author code was executed.\par
The root read all 62 primary pages and synthesized three native bounded checks of extensions, spin kernels, and nonspin/bordism calculations. These auxiliaries are not independent acceptance reviewers and are not attributed to another model. The external controller’s Personal-account/Chat/GPT-6/Pro UI observation remains that controller’s observation; the backend was not independently verified.\par
The principal modified-surgery statement, its good-normal-structure condition, and selected bordism calculations were inspected in [Retained latex: \textbackslash{}cite[Theorems 5, 6, 8 and Section 5]\{Kreck\}]. The original classification, embedding, pseudo-isotopy and surgery literature, Crowley’s integral extension theorem, and the Borel cohomology input were not all reproved. This boundary is material to any assertion that the entire geometric classification has been independently certified. The decisive elementary contradictions above are recorded separately from those dependencies.\par
9 M. Kreck and Y. Su, Mapping class group of manifolds which look like 3-dimensional complete intersections, arXiv:2009.08054v3 (11 June 2024), inspected author manuscript. M. Kreck, On the classification of 1-connected 7-manifolds with torsion free second homology, arXiv:1805.02391v1 (7 May 2018), selected pp. 6–7, 11–14, 20–22. D. Benson, C. Campagnolo, A. Ranicki and C. Rovi, Cohomology of symplectic groups and Meyer’s signature theorem, arXiv:1710.04851v2, selected pp. 14–15. S. Galatius and O. Randal-Williams, Moduli spaces of manifolds: a user’s guide, arXiv:1811.08151v2, Section 5.3. M. Krannich, Mapping class groups of highly connected \((4k+2)\)-manifolds, arXiv:1902.10097v2, Appendix A. P. Teichner, On the signature of four-manifolds with universal covering spin, Math. Ann. 295 (1993), 745–759; selected pp. 750–751, author-hosted published PDF.\par
\end{document}
